\documentclass[a4paper]{article}

% Basic packages
% Español (México) con XeLaTeX: fontspec para fuentes Unicode, polyglossia
% para la división silábica y los nombres del documento en la variante mexicana.
\usepackage{fontspec}
\usepackage{polyglossia}
\setdefaultlanguage[variant=mexican]{spanish}
% Los nombres chinos van como «pinyin (original)»: xeCJK da a los caracteres
% Han su propia fuente; sin ella XeLaTeX los omite con un aviso.
% FandolSong no cubre algunos caracteres de nombres (U+73FA); WenQuanYi Zen Hei sí.
\usepackage[AutoFallBack=true]{xeCJK}
\setCJKmainfont{FandolSong-Regular.otf}
\setCJKfallbackfamilyfont{\CJKrmdefault}{WenQuanYi Zen Hei}
% polyglossia no traduce los nombres de listings.
\AtBeginDocument{\providecommand{\lstlistingname}{}\renewcommand{\lstlistingname}{Listado}%
  \providecommand{\lstlistlistingname}{}\renewcommand{\lstlistlistingname}{Índice de listados}}
\usepackage{amsmath, amssymb}
\usepackage{graphicx}
\usepackage[margin=2.5cm]{geometry}
\usepackage[most]{tcolorbox}
\usepackage{etoolbox}
\usepackage{listings}
\usepackage{hyperref}
\usepackage{booktabs}
\usepackage{subcaption}
\usepackage{float}

% Highlight boxes
\newtcolorbox{knowledgebox}[1]{
    enhanced, colback=blue!5!white, colframe=blue!75!black, colbacktitle=blue!75!black,
    coltitle=white, fonttitle=\bfseries, title={#1}, attach boxed title to top left={yshift=-2mm, xshift=2mm},
    boxrule=1pt, sharp corners
}

\newtcolorbox{importantbox}[1]{
    enhanced, colback=yellow!10!white, colframe=yellow!80!black, colbacktitle=yellow!80!black,
    coltitle=black, fonttitle=\bfseries, title={#1}, sharp corners
}

\newtcolorbox{warningbox}[1]{
    enhanced, colback=red!5!white, colframe=red!75!black, colbacktitle=red!75!black,
    coltitle=white, fonttitle=\bfseries, title={#1}, sharp corners
}

% Listings
\lstset{
    language=Python,
    basicstyle=\ttfamily\small,
    keywordstyle=\color{blue},
    stringstyle=\color{red!60!black},
    commentstyle=\color{green!60!black},
    breaklines=true,
    frame=single,
    numbers=left,
    numberstyle=\tiny\color{gray},
    captionpos=b,
    extendedchars=true
}

% Front-page metadata
\newcommand{\notetitle}{{[}LLM Agents F24{]} LLM Reasoning — Denny Zhou}
\newcommand{\noteauthors}{Elaboradas a partir de materiales públicos del curso}
\newcommand{\notedate}{9 de septiembre de 2024}
\newcommand{\videochannel}{Berkeley RDI}
\newcommand{\videopublishdate}{9 de septiembre de 2024}
\newcommand{\videoduration}{aproximadamente 60 minutos}
\newcommand{\videourl}{https://www.youtube.com/watch?v=QL-FS_Zcmyo}
\newcommand{\videocoverpath}{cover.jpg}
\newcommand{\repourl}{https://github.com/hqhq1025/ai-course-notes}


% Footer with repo info
\usepackage{fancyhdr}
\pagestyle{fancy}
\fancyhf{}
\fancyfoot[C]{\thepage}
\fancyfoot[R]{\footnotesize\color{gray}\href{\repourl}{AI Course Notes}}
\renewcommand{\headrulewidth}{0pt}
\renewcommand{\footrulewidth}{0pt}

\begin{document}

\begin{titlepage}
\centering
{\Large Notas del curso\par}
\vspace{1.2cm}
{\huge\bfseries \notetitle\par}
\vspace{0.8cm}
{\large \noteauthors\par}
\vspace{0.3cm}
{\large \notedate\par}
\vspace{1.2cm}

\ifdefempty{\videocoverpath}{
{\small Al generar las notas, indique la ruta local de la portada del video.\par}
}{
\includegraphics[width=0.82\textwidth,height=0.45\textheight,keepaspectratio]{\videocoverpath}\par
}

\vfill
\begin{tcolorbox}[width=0.9\textwidth, colback=black!2!white, colframe=black!60, sharp corners]
\textbf{Autor o canal del video}: \videochannel\par
\textbf{Fecha de publicación}: \videopublishdate\par
\textbf{Duración del video}: \videoduration\par
\ifdefempty{\videourl}{}{
\textbf{Enlace del video}: \href{\videourl}{\nolinkurl{\videourl}}\par
}
\par
\textbf{Página del proyecto}: \href{\repourl}{\nolinkurl{\repourl}}\par
\end{tcolorbox}
\end{titlepage}

\tableofcontents
\newpage

%% --- Inicio del contenido --- %%

\section{Introducción: por qué la capacidad de razonamiento es clave para la IA}

Denny Zhou es científico investigador en Google DeepMind, dedicado desde hace mucho tiempo a la investigación de la capacidad de razonamiento de los modelos de lenguaje grandes. Es uno de los contribuyentes centrales de trabajos histĂłricos como Chain-of-Thought Prompting y Self-Consistency. Esta conferencia parte de una pregunta fundamental —\textbf{¿qué es realmente la inteligencia?}— para desarrollar el hilo central del razonamiento en los LLM.

\begin{knowledgebox}{El cambio de paradigma del aprendizaje automático al razonamiento}
El aprendizaje automático tradicional (few-shot learning, active learning, meta-learning, etc.) persigue la eficiencia de datos (data efficiency), pero en la práctica estos métodos no logran realmente «aprender a partir de pocas muestras». Denny Zhou sostiene que la razón fundamental por la que los seres humanos pueden aprender a partir de pocas muestras no son las regularidades estadísticas, sino la \textbf{capacidad de razonamiento (reasoning)}. Esta intuición impulsó su transición del aprendizaje automático hacia la investigación del razonamiento en los LLM.
\end{knowledgebox}

\subsection{Last Letter Concatenation: un ejemplo simple pero profundo}

Denny Zhou utiliza la tarea de Last Letter Concatenation para revelar la enorme diferencia entre los métodos tradicionales y los métodos de razonamiento:

\begin{itemize}
    \item \textbf{Definición de la tarea}: dado el nombre de una persona (por ejemplo, ``Elon Musk''), producir la concatenación de la última letra del nombre y la última letra del apellido (``n'' + ``k'' = ``nk'').
    \item \textbf{Método tradicional de aprendizaje automático}: entrenado con un Transformer, requiere miles de muestras etiquetadas y alcanza una exactitud de aproximadamente 85--90\%.
    \item \textbf{Few-shot Prompting}: al darle directamente al LLM unos cuantos ejemplos, este sigue cometiendo errores con facilidad (por ejemplo, para ``Barack Obama'' produce ``ck'' en lugar de ``ka'').
    \item \textbf{Chain-of-Thought Prompting}: al agregar pasos de razonamiento dentro de los ejemplos (``The last letter of Elon is n, the last letter of Musk is k, so nk''), basta con un solo ejemplo para alcanzar una exactitud de 100\%.
\end{itemize}

\begin{importantbox}{Intuición central}
Para una tarea que resulta simple para los seres humanos, si un método necesita una cantidad masiva de datos etiquetados para aprenderla, entonces es difícil considerarlo verdadera «inteligencia». La capacidad de razonamiento permite que los LLM generalicen a partir de muy pocas muestras, algo que los métodos tradicionales de aprendizaje automático no pueden igualar.
\end{importantbox}

\subsection{Resumen de la sección}

La capacidad de razonamiento es la clave para cerrar la brecha entre el «aprendizaje impulsado por datos» y la «inteligencia real». Al generar pasos intermedios al estilo Chain-of-Thought, los LLM logran un salto cualitativo: pasar de necesitar grandes cantidades de datos de entrenamiento a requerir solo unas cuantas muestras, o incluso ninguna.
\section{Pasos intermedios (Intermediate Steps): el mecanismo central del razonamiento}

\subsection{Trayectoria histórica: del entrenamiento al prompting}

Denny Zhou trazó el desarrollo de la idea de los pasos intermedios:

\begin{enumerate}
    \item \textbf{2017, Ling et al.}: propusieron en su artículo resolver problemas matemáticos mediante pasos de razonamiento en lenguaje natural (rationale), entrenando modelos seq2seq para generar los pasos intermedios desde cero y llegar así a la respuesta. Denny Zhou lo calificó de «como un viajero en el tiempo».
    \item \textbf{2021, OpenAI GSM8K}: dando continuidad a la idea de 2017, crearon un conjunto de datos matemático con pasos intermedios para hacer fine-tuning de GPT-3, lo que mejoró de manera considerable la capacidad de razonamiento matemático.
    \item \textbf{2021, Google Brain Scratchpad}: descubrieron de forma independiente una idea similar, pero la aplicaron en el dominio de la síntesis de programas (program synthesis) usando pasos intermedios simbólicos.
    \item \textbf{2022, Chain-of-Thought Prompting (Wei et al.)}: evaluaron de manera extensa el efecto de los pasos intermedios en la etapa de prompting, mostrando mejoras notables en prácticamente todas las tareas de PLN.
\end{enumerate}

\begin{importantbox}{Los pasos intermedios son la clave, no el método concreto}
Ya sea mediante training, fine-tuning o prompting, lo que en realidad funciona son los \textbf{pasos intermedios (intermediate steps)} en sí mismos. Cuando se proporcionan ejemplos que incluyen pasos intermedios, el LLM genera respuestas que también los incluyen; esto es una consecuencia natural de que el LLM, como modelo probabilístico, imita el patrón de la entrada.
\end{importantbox}

\subsection{Fundamento teórico: por qué funcionan los pasos intermedios}

Denny Zhou citó un trabajo teórico de 2024 realizado en colaboración con Branden Srouji, de Stanford:

\begin{itemize}
    \item \textbf{Transformer con pasos intermedios}: basta con que la profundidad supere una \textbf{constante} (independiente de la longitud de la entrada) para poder resolver cualquier problema inherentemente serial (inherently serial).
    \item \textbf{Transformer que produce la respuesta directamente}: o bien requiere una profundidad enorme, o bien es imposible de resolver.
\end{itemize}

\begin{warningbox}{Implicación práctica}
Si un modelo no logra resolver cierto problema, una estrategia eficaz es guiarlo para que genere más pasos intermedios. Además, se pueden invocar herramientas externas para apoyar la generación de los pasos intermedios; esta es precisamente una de las ideas centrales del marco de los LLM Agent.
\end{warningbox}

\subsection{Resumen de la sección}

La idea de los pasos intermedios pasó por una evolución que fue del entrenamiento al fine-tuning y de ahí al prompting, pero su esencia se mantuvo invariable: hacer que el modelo derive la respuesta paso a paso en lugar de saltar directamente a ella. En términos teóricos, los pasos intermedios dotan a un Transformer de profundidad finita de la capacidad de resolver cualquier cómputo serial.

\section{Estrategias de razonamiento: de Few-shot a Zero-shot}

\subsection{Least-to-Most Prompting: descomposición de tareas complejas}

Inspirado en la estrategia de descomposición de la obra clásica de P\'{o}lya \textit{How to Solve It}, el equipo de Denny Zhou propuso Least-to-Most Prompting:

\begin{itemize}
    \item \textbf{Idea central}: primero se descompone el problema complejo en una serie de subproblemas, luego se resuelven de manera progresiva a partir del subproblema más simple, y la solución de cada subproblema sirve como contexto para el siguiente.
    \item \textbf{Tarea SCAN}: punto de referencia de generalización composicional (compositional generalization), donde se alcanza una exactitud de 99.7\% usando solo 1\% de los ejemplos.
    \item \textbf{Tarea Text-to-Code}: usando Dynamic Least-to-Most Prompting, con solo 1\% de los datos se supera ampliamente a arquitecturas especializadas entrenadas con todos los datos de entrenamiento.
\end{itemize}

\begin{knowledgebox}{Generalización composicional (Compositional Generalization)}
Se refiere al escenario en el que las muestras de prueba son más complejas que las muestras de entrenamiento o de los prompts. Por ejemplo, durante el entrenamiento solo se han visto fragmentos de código cortos, pero en la prueba se requiere generar código más largo. Least-to-Most Prompting resuelve este problema de manera natural mediante su estrategia de descomposición.
\end{knowledgebox}

\subsection{Zero-shot CoT: ``Let's think step by step''}

\begin{itemize}
    \item Sin necesidad de ningún ejemplo, basta con añadir la frase ``Let's think step by step'' después de la pregunta para activar el razonamiento paso a paso del LLM.
    \item El desempeño suele ser inferior al de few-shot CoT, pero tiene la ventaja de no tener costo alguno.
\end{itemize}

\subsection{LLMs as Analogical Reasoners}

Inspirado en la idea de razonamiento analógico de P\'{o}lya:

\begin{itemize}
    \item \textbf{Método}: ante un problema nuevo, se hace que el LLM genere por sí mismo problemas relacionados junto con sus soluciones, y luego se utilizan esos ejemplos autogenerados para resolver el problema original.
    \item \textbf{Ventaja}: es notablemente mejor que ``Let's think step by step'', e incluso supera a los ejemplos few-shot diseñados manualmente, porque el modelo genera ejemplos relacionados hechos a la medida de cada problema.
    \item \textbf{Relación con la generación aumentada por recuperación}: puede extenderse además a la recuperación de problemas y conocimientos relacionados desde la red, lo que permite lograr escalabilidad.
\end{itemize}

\subsection{Chain-of-Thought Reasoning without Prompting}

Un hallazgo sorprendente: \textbf{sin necesidad de ningún prompt}, es posible activar el razonamiento únicamente mediante una estrategia de decodificación especial:

\begin{itemize}
    \item En el primer paso de decodificación, en lugar de tomar el token de mayor probabilidad, se exploran los $k$ tokens candidatos principales.
    \item Para cada token candidato se continúa con decodificación voraz, generando una respuesta completa.
    \item Se elige la ruta que contiene pasos de razonamiento y cuya respuesta final tiene la mayor confianza.
\end{itemize}

\begin{importantbox}{Observación clave}
El LLM ya «sabe» internamente cómo razonar. Cuando aparece una ruta de razonamiento (como ``Nicolas Cage was born in 1964, and 1964 is an even year''), la probabilidad de la respuesta final salta de un valor bajo a 98\%. Los pasos de razonamiento aumentan de manera significativa, a nivel de probabilidad, la confianza del modelo en la respuesta final.
\end{importantbox}

\subsection{Resumen de la sección}

De few-shot CoT a zero-shot CoT, y luego a la estrategia de decodificación especial que no requiere prompt alguno, las formas de activar la capacidad de razonamiento del LLM son cada vez más flexibles. La regularidad central es que lo determinante son los pasos de razonamiento en sí mismos, y no la forma de activarlos. El razonamiento analógico muestra además que el LLM puede generar de manera autónoma ejemplos relacionados para apoyar su razonamiento.
\section{Self-Consistency: del muestreo a la votación}

\subsection{Motivación: la naturaleza probabilística de los LLM}

Denny Zhou parte de primeros principios para derivar la validez de Self-Consistency:

\begin{itemize}
    \item Al decodificar, el LLM optimiza: $\arg\max P(\text{reasoning path}, \text{answer} \mid \text{problem})$
    \item Pero lo que en realidad necesitamos es: $\arg\max P(\text{answer} \mid \text{problem})$
    \item La relación entre ambos: $P(\text{answer} \mid \text{problem}) = \sum_{\text{path}} P(\text{answer}, \text{path} \mid \text{problem})$
    \item Es decir, hay que sumar (marginalizar) sobre todas las rutas de razonamiento posibles, en lugar de tomar solo la única ruta óptima.
\end{itemize}

\subsection{Método y resultados}

\begin{importantbox}{Método de Self-Consistency}
\begin{enumerate}
    \item Muestrear el mismo problema varias veces (con temperatura distinta de cero) para obtener varias rutas de razonamiento y sus respuestas correspondientes.
    \item Tomar la \textbf{respuesta con mayor frecuencia de aparición} como predicción final (votación mayoritaria / majority voting).
    \item Nota: lo que se vota es la \textbf{respuesta final}, no la ruta de razonamiento; la ruta de razonamiento es una variable latente (latent variable).
\end{enumerate}
\end{importantbox}

Self-Consistency trajo una enorme mejora de desempeño y en su momento superó ampliamente el estado del arte en todos los puntos de referencia. Más importante aún, cuando la consistencia supera el 80\%, la exactitud se acerca al 100\%, lo cual ofrece un indicador confiable de confianza.

\begin{knowledgebox}{Extensión con Universal Self-Consistency}
Para respuestas de forma libre (free-form), el propio LLM puede juzgar qué respuestas son semánticamente equivalentes y así encontrar el grupo de respuestas más común. Por ejemplo, para una pregunta como «qué países tienen un consumo per cápita de café menor que el de México», distintas respuestas pueden tener una redacción diferente pero un contenido consistente.
\end{knowledgebox}

\subsection{Resumen de la sección}

Self-Consistency eleva el razonamiento del LLM de una única generación al nivel de la inferencia estadística. Su idea central proviene de la inferencia marginal (marginal inference) de los modelos gráficos probabilísticos: es simple pero sumamente efectiva.

\section{Limitaciones del razonamiento}

\subsection{Los LLM se distraen fácilmente con contexto irrelevante}

\begin{itemize}
    \item Al insertar información irrelevante en un problema de matemáticas (como ``Mario's moustache costs \$10''), el LLM se deja engañar.
    \item Agregar la instrucción ``ignore irrelevant context'' mitiga el problema parcialmente, pero su efecto es limitado cuando aumenta la cantidad de contenido irrelevante.
    \item Incluso la acumulación masiva de oraciones irrelevantes sencillas (``the sky is blue'') provoca una caída significativa del desempeño.
\end{itemize}

\begin{warningbox}{Interferencia del contexto}
El LLM, como modelo probabilístico, incorpora todos los tokens de entrada al cálculo de atención. La información irrelevante no solo desperdicia la ventana de contexto, sino que interfiere de manera sustancial con el proceso de razonamiento. Esto exige especial atención en las aplicaciones reales de Agentes: el ruido de los documentos que introduce la generación aumentada por recuperación (RAG) puede, por el contrario, perjudicar la calidad del razonamiento.
\end{warningbox}

\subsection{El LLM no puede autocorregir su razonamiento}

Los experimentos del equipo de Denny Zhou muestran que:

\begin{itemize}
    \item Cuando se le pide al LLM que revise y corrija su propia respuesta, puede corregir una respuesta incorrecta, pero también puede convertir una respuesta correcta en incorrecta.
    \item En bases de referencia como GSM8K y CommonsenseQA, los métodos de self-correction \textbf{no produjeron ninguna mejora neta}; al contrario, empeoraron el desempeño.
    \item Las mejoras de self-correction reportadas en la literatura suelen usar un Oracle (es decir, solo se le pide al modelo que corrija cuando la respuesta ya es incorrecta), pero el modelo mismo no puede determinar si su respuesta es correcta.
\end{itemize}

\begin{importantbox}{La retroalimentación externa es condición para la autocorrección}
La autocorrección del LLM requiere retroalimentación externa de un Oracle (como las pruebas unitarias en tareas de código). El trabajo de Self-Debug ofrece este Oracle de manera natural mediante pruebas unitarias. Depender únicamente del juicio del propio modelo para corregirse no es confiable.
\end{importantbox}

\subsection{El Multi-Agent Debate no supera a Self-Consistency}

\begin{itemize}
    \item Varios Agentes debaten entre sí y llegan a un consenso (multi-agent debate), generando en total $n$ respuestas.
    \item Sin embargo, aplicar simplemente una votación de Self-Consistency sobre $n$ muestras independientes supera de manera consistente al multi-agent debate.
    \item La interacción de información que introduce el proceso de debate no aporta ningún beneficio adicional.
\end{itemize}

\subsection{El orden de las premisas (Premise Order) afecta el razonamiento}

\begin{itemize}
    \item \textbf{Experimento con GSM8K}: al reordenar únicamente las oraciones del problema (sin cambiar el significado), la exactitud cae alrededor de 10 puntos porcentuales.
    \item \textbf{Experimento de razonamiento lógico}: usando razonamiento puramente lógico con reglas de símbolos aleatorios, con solo alterar el orden de las reglas relevantes, todos los modelos de frontera muestran una caída de más de 30 puntos porcentuales en la exactitud.
    \item \textbf{Causa raíz}: el LLM solo puede procesar la entrada de manera secuencial y no puede, como los humanos, saltar libremente entre premisas ni retroceder sobre ellas.
\end{itemize}

\begin{warningbox}{Efecto del orden de las premisas}
El razonamiento del LLM depende en gran medida del orden en que se presentan las premisas. Al diseñar sistemas de Agentes, el orden en que se organiza la información de entrada debe alinearse con el orden que requiere el razonamiento; de lo contrario, se produce una pérdida significativa de desempeño.
\end{warningbox}

\subsection{Resumen de la sección}

El razonamiento de los LLM actuales presenta tres limitaciones principales: se distrae fácilmente con contexto irrelevante, no puede autocorregirse de manera confiable y es sensible al orden de las premisas. Estas limitaciones ofrecen una guía importante para el diseño de sistemas de Agentes: es necesario compensar estas debilidades mediante herramientas externas y flujos de trabajo cuidadosamente diseñados.
\section{Síntesis y ampliación}

\subsection{Repaso de los puntos centrales}

\begin{enumerate}
    \item Los \textbf{pasos intermedios} son el mecanismo central para mejorar la capacidad de razonamiento de los LLM, ya sea mediante training, fine-tuning o prompting.
    \item \textbf{Self-Consistency} mejora de manera significativa la exactitud del razonamiento desde una perspectiva probabilística, mediante muestreos múltiples y votación mayoritaria.
    \item Las \textbf{estrategias de razonamiento} van desde few-shot CoT hasta zero-shot CoT, y de ahí a analogical reasoning y decodificación sin prompt, con formas de activación cada vez más flexibles.
    \item Las \textbf{limitaciones} —interferencia del contexto, fallas de autocorrección, sensibilidad al orden de las premisas— señalan las direcciones de mejora futuras.
\end{enumerate}

\subsection{Implicaciones para los LLM Agent}

\begin{itemize}
    \item \textbf{Llamadas a herramientas y retroalimentación externa}: un Agent puede proporcionar un Oracle externo llamando a ejecutores de código, motores de búsqueda y otras herramientas, lo cual compensa la incapacidad de los LLM para autocorregirse.
    \item \textbf{La generación aumentada por recuperación requiere cautela}: los documentos que introduce el RAG pueden contener información irrelevante que, en cambio, interfiere con el razonamiento; se necesita una recuperación y un filtrado precisos.
    \item \textbf{Organización de la información}: el orden en que se presenta la información en el flujo de trabajo de un Agent debe alinearse con la lógica del razonamiento.
    \item \textbf{Escalamiento del razonamiento}: Self-Consistency y Analogical Reasoning pueden servir como medios de escalamiento del cómputo en tiempo de inferencia.
\end{itemize}

\subsection{Lecturas adicionales}

\begin{itemize}
    \item Wei et al., ``Chain-of-Thought Prompting Elicits Reasoning in Large Language Models,'' NeurIPS 2022.
    \item Wang et al., ``Self-Consistency Improves Chain of Thought Reasoning in Language Models,'' ICLR 2023.
    \item Zhou et al., ``Least-to-Most Prompting Enables Complex Reasoning in Large Language Models,'' ICLR 2023.
    \item Kojima et al., ``Large Language Models are Zero-Shot Reasoners,'' NeurIPS 2022.
    \item Yasunaga et al., ``Large Language Models as Analogical Reasoners,'' ICLR 2024.
    \item Huang et al., ``Large Language Models Cannot Self-Correct Reasoning Yet,'' ICLR 2024.
    \item P\'{o}lya, ``How to Solve It,'' Princeton University Press, 1945.
\end{itemize}

%% --- Fin del contenido --- %%

\end{document}
